
% !TeX spellcheck = pt_BR
% !TEX encoding = UTF-8 Unicode

\documentclass[12pt,a4paper]{article}

% TROCAR PELO SEU PRÓPRIO NOME:
\newcommand{\myauthor}{Leonardo~T.~Rolla}

% TROCAR PELO SEU PRÓPRIO REPOSITÓRIO:
%\newcommand{\mysourceurl}{\url{https://github/xxxxxxxxxx}}
\newcommand{\mysourceurl}{\url{https://github.com/rea-stat/Testes-de-Hipotese} e \url{https://www.overleaf.com/read/dkbvxdjftbbd}}

\usepackage[brazil]{babel}

\usepackage[utf8]{inputenc}
\usepackage[T1]{fontenc}
\usepackage{lmodern}

\usepackage{amsfonts}
\usepackage{amsmath}
\usepackage{amssymb}
\usepackage{amsthm}
%\usepackage{chngcntr}
\usepackage{dsfont}
\usepackage[shortlabels]{enumitem}
\usepackage{geometry}
\usepackage{graphicx}
\usepackage{hyperref}
\usepackage{float}
\usepackage{icomma}
\usepackage{mathrsfs}  
\usepackage{microtype}
\usepackage{verbatim}
\usepackage[dvipsnames]{xcolor}

\usepackage{upquote}

\theoremstyle{plain}
\newtheorem{theorem}[equation]{Teorema}
\newtheorem{proposition}[equation]{Proposição}
\newtheorem{lemma}[equation]{Lema}
\newtheorem{corollary}[equation]{Corolário}
\newtheorem{axiom}[equation]{Axioma}
\theoremstyle{definition}
\newtheorem{definition}[equation]{Definição}
\newtheorem*{definition*}{Definição}
\theoremstyle{remark}
\newtheorem{examplex}[equation]{Exemplo}
\newtheorem{counterexamplex}[equation]{Contra-exemplo}
\newtheorem{exercisex}[equation]{Exercício}
\newtheorem{comentario}[equation]{Comentário}
\newtheorem{remarkx}[equation]{Observação}
\newtheorem*{remark*}{Observação}

\newenvironment{example}{\pushQED{\qed}\renewcommand{\qedsymbol}{\scriptsize$\triangle$}\examplex}{\popQED\endexamplex}
\newenvironment{counterexample}{\pushQED{\qed}\renewcommand{\qedsymbol}{\scriptsize$\triangle$}\counterexamplex}{\popQED\endcounterexamplex}
\newenvironment{exercise}{\pushQED{\qed}\renewcommand{\qedsymbol}{\scriptsize$\triangle$}\exercisex}{\popQED\endexercisex}
\newenvironment{remark}{\pushQED{\qed}\renewcommand{\qedsymbol}{\scriptsize$\triangle$}\remarkx}{\popQED\endremarkx}

\geometry{margin=25mm,centering,a4paper}

\setlist{topsep=0em, partopsep=.5em, itemsep=0pt, noitemsep}
\setlist[enumerate,1]{label=\rm(\alph*)}
\setlist[enumerate,2]{label=(\roman*)}

\hypersetup{
hidelinks,
colorlinks,
urlcolor=[rgb]{0,0,.7},
linkcolor=[rgb]{.4,0,0},
citecolor=[rgb]{0,.3,0},
pdfstartview=FitH,
pdfpagemode=UseNone
}

\renewcommand{\baselinestretch}{1.1}
\setlength{\parindent}{0pt}
\setlength{\parskip}{.4em}

\renewcommand*{\url}[1]{\href{{#1}}{\sf\footnotesize\detokenize{#1}}}

\graphicspath{{figures/}}

\usepackage[main,input,include]{embedall}
%Use \embedfile{xxx.pdf} to manually embed figures!

\begin{document}


\title{Testes de Hipótese}
\author{\myauthor}
\maketitle

\thispagestyle{empty}


\cleardoublepage


%\vspace*{\fill}

\noindent
\copyright\
2023--{\the\year}
Departamento de Estatística, IME-USP.
\\

\vspace{1em}

\hfil
\href{https://creativecommons.org/licenses/by-sa/4.0/deed.pt-br}{\includegraphics[height=4.2em]{by-sa.pdf}}
\embedfile{figures/by-sa.pdf}

\vspace{1em}

\noindent
Todos os direitos reservados.
Permitido o uso nos termos licença Creative Commons Atribuição-CompartilhaIgual 4.0 Internacional.

\vspace{1em}

\textbf{Reutilização deste material}
\\
\noindent
Você pode remixar, transformar, e criar a partir do material para qualquer fim, mesmo que comercial.
Nesse caso, tem de distribuir as suas contribuições sob a mesma licença que o original.
Você não pode aplicar termos jurídicos ou medidas de caráter tecnológico que restrinjam legalmente outros de fazerem algo que a licença permita.

\vspace{1em}

\textbf{Atribuição}
\\
\noindent
Este material foi produzido originalmente por Leonardo~T.~Rolla (IME-USP), como parte do projeto \emph{REA-Stat -- Recursos Educacionais Abertos de Estatística}
(\url{https://rea-stat.github.io/}).

\vspace{1em}

\textbf{Código-fonte}
\\
\noindent
O código-fonte deste material está disponível em:
\\
\mysourceurl


\vspace{1em}

\textbf{Aviso legal}
\\
\noindent
As pessoas e instituições aqui mencionadas não endossam a qualidade deste material e as opiniões nele contido, nem explícita nem implicitamente.

\noindent
Qualquer erro contido neste material é responsabilidade de: \myauthor.

\vspace{1em}

\hfill
\today

\vfill
\vfill

\thispagestyle{empty}


\cleardoublepage


\section{Testes de hipótese para uma proporção}

Começamos com alguns exemplos para ir introduzindo os principais conceitos envolvidos.

\begin{example}
\label{exp:facil}
Numa cidade do interior, todos os administradores de estabelecimento com piscinas de bolinhas para crianças compram as bolinhas da Sra.~Mariana, uma mulher bem sucedida que há vários anos está à frente de uma distribuidora desses produtos.
Das bolinhas vendidas pela distribuidora, $70\%$ são de cores bonitas e vivas (azul, verde, preto, laranja, vermelho, fúcsia, violeta, branco, etc.), e as restantes são um pouco menos vivas e bonitas (manchadas, acinzentadas, etc.).
Pelo menos é isso o que Sra.~Mariana sempre disse, e os empresários da região sempre acreditaram.

Um dia, o Sr.~Rubens descobriu que havia uma outra indústria que fabricava um produto muito parecido, porém com qualidade um pouco inferior, tendo apenas $60\%$ das bolas em cores vivas.
Ele começou a suspeitar que a Sra.~Mariana estava distribuindo o produto dessa nova indústria.
Depois de muito hesitar, o Sr.~Rubens eventualmente decidiu compartilhar sua preocupação com os demais empresários da cidade.

Para colocar à prova a premissa de que a Sra.~Mariana vinha cumprindo com os parâmetros prometidos, a comunidade decidiu que deveriam tomar $n=1.000$ amostras ao acaso de um grande lote de bolinhas, e observar quantas seriam de cores vivas.
Assim, definiu-se a estatística
\[
\hat{p} = \frac{\text{número de bolas de cores vivas}}{n}
\]
e decidiu-se pelo seguinte procedimento.
Se $\hat{p} \leq 0,65$, confirmariam-se as suspeitas do Sr.~Rubens, rejeitando a premissa de que $70\%$ das bolas eram de cores vivas.
Do contrário, se $\hat{p}>0,65$, a reputação da Sra.~Mariana continuaria inabalada.

A Sra.~Mariana fez seus cálculos e concluiu que a probabilidade de que fosse cometida uma injustiça contra ela era menor do que $0,1\%$.
Quer dizer, supondo que $p=0,7$, ela calculou $P(\hat{p} \leq 0,65) < 0,001$.\footnote{Os cálculos feitos aqui usam a aproximação de $\hat{p}$ por uma normal. Os próximos exemplos dão mais detalhes desse método, e os cálculos feitos em R são exibidos ao final desta seção.}
Assim sendo, a Sra.~Mariana não se opôs à realização do experimento.
Por sua vez, o Sr.~Rubens fez outros cálculos e observou que, caso sua suspeita fosse verdadeira, as chances de que não fossem confirmadas também seriam menores que $0,1\%$.
Quer dizer, supondo que $p=0,6$, ele calculou $P(\hat{p} > 0,65) < 0,001$.
Assim sendo, o Sr.~Rubens tampouco se opôs à realização do experimento.
Dessa forma, a situação foi resolvida sem grandes controvérsias: bastaria fazer a amostra e decidir a partir dela, com um critério bem claro.
\end{example}

O exemplo acima apresentou uma situação muito típica: acreditava-se em uma determinada hipótese, até que eventualmente decidiu-se que ela deveria ser posta à prova.
A hipótese aceita até então recebe o nome de \emph{hipótese nula}, que agora passaria a ser desafiada por uma \emph{hipótese alternativa}.
O experimento planejado e o critério de decisão adotados foram muito claros.
Nesse exemplo, o critério de decisão não gerou nenhum dilema importante.
Porém, isso só foi possível graças a um grande dispêndio de recursos, o que muitas vezes é impossível, inviável ou indesejado.
Numa situação um pouco mais realística, começamos a perceber alguns conflitos e ter que tomar decisões em face de dilemas consideráveis.

\begin{example}
\label{exp:dificil}
Continuando o exemplo anterior, a comunidade eventualmente percebeu que não tinha os recursos necessários para amostrar e analisar $1.000$ bolinhas, e pensou em fazer o mesmo experimento, porém com apenas $n=100$ amostras.

Ao saber disso, a Sra.~Mariana imediatamente fez as contas, obtendo $P(\hat{p} \leq 0,65) = 0,14$. Como esperado, ela protestou, pois considerou injusto ficar à mercê de um procedimento aleatório com chances consideráveis de manchar sua reputação injustamente.

Algumas pessoas imediatamente fizeram um outro cálculo e descobriram que, para que as chances de cometer uma injustiça continuassem menores que $0,001$, teriam que adotar o nível de corte como sendo $0,558$ ao invés de $0,65$.
Mas isso lhes pareceu inadequado, pois, supondo que $p=0,6$, teríamos $P(\hat{p} > 0,558) = 0,80$.
Ou seja, mesmo se as suspeitas do Sr.~Rubens estivessem corretas, haveria $80\%$ de probabilidade de que elas não fossem confirmadas.
Isso comprometeria muito as chances de que o teste tivesse alguma utilidade.

Chamaram então o filho do Sr.~Januário, que na época estudava numa universidade na capital e já havia até cursado algumas disciplinas de Estatística.
O filho do Sr.~Januário disse que teriam que tomar como critério $\hat{p} \leq 0,625$ para rejeitar a premissa de honestidade da Sra.~Mariana.
Ela imediatamente fez as contas, e disse que isso ainda lhe parecia um pouco injusto, pois haveria uma probabilidade $0,05$ de que a premissa fosse rejeitada incorretamente.
O filho do Sr.~Januário disse que não tinha muito o que fazer, que esse procedimento é inerentemente aleatório, e que $\alpha=0,05$ lhe parecia apropriado neste caso.
Os empresários não haviam cursado Estatística na universidade mas eram bastante espertos, e indagaram qual a justificativa desse valor de $\alpha=0,05$.
O filho do Sr.~Januário não soube explicar, disse que está bem fazer desse jeito e que é assim mesmo que se faz.
Eventualmente, a Sra.~Mariana se resignou e consentiu.

A essa altura, já podemos imaginar que quem protestou foi o Sr.~Rubens.
Ele disse que, se suas suspeitas fossem bem fundadas, ainda assim haveria probabilidade $31\%$ de que o teste não as confirmasse.
O filho do Sr.~Januário disse que não podia fazer nada a respeito, explicando que esse valor, denotado por $\beta$, era menos importante, e que o procedimento deveria, em primeiro lugar, reduzir as chances de que uma injustiça fosse cometida contra a reputação da Sra.~Mariana.
Como consolo, o filho do Sr.~Januário disse que, apesar de que o valor de $\alpha$ já tinha sido fixado e não deveria ser questionado, seria possível diminuir o valor de $\beta$ tomando-se uma amostra maior do que $n=100$.
Ao final, ele sugeriu ao Sr.~Rubens não se focar muito no valor de $\beta$, e sim no seu complementar, ou seja, os $69\%$ de chance que ele terá de ver suas suspeitas confirmadas, caso sejam verdadeiras.

Antes de encerrar a discussão e proceder à amostragem, o filho do Sr.~Januário achou importantíssimo enfatizar que tipo de interpretação poderíamos fazer do resultado desse teste.
Ele esclareceu que, após a realização da amostra, haveria duas possibilidades bem distintas em sua natureza.
Na primeira possibilidade, de $\hat{p} \leq 0,625$, ficaria estabelecido que \emph{houve suficiente evidência estatística} para rejeitar a premissa de que $70\%$ das bolinhas da Sr.~Mariana eram de cores vivas.
Na segunda possibilidade, de $\hat{p} > 0,625$, ficaria estabelecido que \emph{não houve suficiente evidência estatística} para rejeitar a premissa.
Nenhuma das duas possibilidades resultaria na \emph{aceitação} ou \emph{confirmação} da premissa original, muito menos permitiria rejeitar a hipótese alternativa.
Essas ponderações deixaram o Sr.~Rubens bem mais tranquilo.
Os empresários da região se sentiram muito melhor esclarecidos e agradeceram pela explicação.
\end{example}

O desenvolvimento acima tornou a estória muito mais interessante, revelando certos conflitos e fazendo transparecer muito mais a assimetria inerente a um teste de hipótese.

Vejamos o contraste entre as hipóteses:

\textbf{Hipótese nula} (``$H_0$''): afirmação sobre o parâmetro $p$, ligada a um valor de referência, um especificação padrão ou histórica. É a hipótese nula que estamos colocando à prova.

\textbf{Hipótese alternativa}  (``$H_1$''): é a hipótese que pode ser considerada aceitável, caso $H_0$ seja rejeitada; afirmação que suspeitamos ser verdadeira sobre a variável em questão.

O resultado do teste se limita a dizer se houve ou não suficiente evidência estatística para rejeitar $H_0$, sem jamais analisar o mérito específico de $H_1$.
Devido a flutuações aleatórias inerentes ao processo de amostragem, a estatística $\hat{p}$ dificilmente será exatamente igual ao seu valor médio previsto por $H_0$, mesmo que essa hipótese seja verdadeira.
Considerando o tamanho típico das flutuações de $\hat{p}$ que esperaríamos observar sob a hipótese nula, diremos que há suficiente evidência estatística para rejeitar tal hipótese quando a discrepância observada for atipicamente grande.
Caso as flutuações não sejam atipicamente grandes para $H_0$, geralmente não poderemos concluir nada sobre $H_0$, $H_1$, ou qualquer hipótese a respeito da variável em questão, e diremos simplesmente que $H_0$ não foi rejeitada.

A partir dessa perspectiva, podemos analisar melhor as preocupações envolvidas num teste de hipótese, e observar as diferenças em natureza e importância que há entre elas.

O \textbf{erro tipo I} consiste em rejeitar a hipótese nula quando ela é verdadeira. Esse erro corresponde a uma \emph{falsa descoberta}. Este é o erro mais grave.
O \textbf{erro tipo II} consiste em não rejeitar a hipótese nula, mesmo ela sendo falsa. Esse erro corresponde à \emph{perda de oportunidade} de se fazer uma descoberta, por falta de suficiente evidência estatística.

Resumimos os possíveis resultados de um teste de hipótese na tabela abaixo:
\begin{table}[H]
\centering
\begin{tabular}{||c|c|c||}
\hline
\hline
Situação & Rejeitamos $H_0$ & Não rejeitamos $H_0$
\\
\hline
$H_0$ é verdadeira & erro tipo I & sem erro
\\
\hline
$H_0$ é falsa & sem erro & erro tipo II
\\
\hline
\hline
\end{tabular}
\end{table}

Antes de realizar o teste, vamos fixar de antemão um parâmetro que corresponderá à probabilidade de que ocorra um erro do tipo I.
Essa probabilidade, denotada pela letra $\alpha$, é chamada de \textbf{nível de significância} do teste.
De forma mais precisa, $\alpha$ é a probabilidade de se rejeitar $H_0$ no caso em que $H_0$ seja verdadeira.
Uma vez fixado o valor de $\alpha$, usamos a hipótese $H_0$ e o número $n$ de amostras para planejar o teste. No exemplo acima, o limiar que separaria a rejeição da não-rejeição foi calculado dessa forma (os cálculos não foram feitos explicitamente, alguma explicação será dada no próximo exemplo).

Observamos que o teste ``protege'' a hipótese nula através da escolha de um parâmetro $\alpha$ pequeno.
Este valor deve ser fixado antes da coleta da amostra.
Dependendo da gravidade do erro de tipo~I, $\alpha$ costuma ser fixado como
$0,05$, $0,01$, ou $0,001$.
Como vimos no exemplo acima, a escolha do $\alpha$ é, de certa forma, arbitrária.
Uma vez fixado o valor de $\alpha$, o ``ônus da prova'' ficará com $H_1$, pois a amostra deverá fornecer evidência estatística suficientemente forte para que o teste seja conclusivo na direção de rejeitar $H_0$.
Isso pode ser muito difícil, sobretudo se $\alpha$ for muito pequeno ou o número $n$ de amostras não for grande o suficiente.
Por outro lado, quanto menor o valor de $\alpha$, mais significante será uma amostra que leve à rejeição de $H_0$. Isso se deve ao fato de que, à medida em que diminuímos o valor $\alpha$, mais a amostra tem que ser incompatível com $H_0$ para levar à rejeição.

Vimos também no final do exemplo algumas limitações com respeito às conclusões que poderemos tirar de um teste de hipótese.
Ao rejeitar a hipótese nula, nunca vamos saber com certeza absoluta se estamos fazendo uma descoberta legítima ou cometendo um erro do tipo I.
Ao não rejeitar a hipótese nula, nunca vamos saber com certeza se estamos deixando de fazer uma descoberta por causa do rigor científico de haver fixado um $\alpha$ muito baixo, ou se a hipótese nula é realmente adequada para modelar a variável em questão.
Em todo caso, o teste nunca servirá para confirmar $H_0$ nem para rejeitar $H_1$.

Também vimos que prefere-se falar de \textbf{poder} do teste, que é a probabilidade de se fazer uma descoberta no caso em que $H_1$ é verdadeira, ao invés do seu complementar $\beta$, que é a probabilidade de erro tipo II.
O poder é representado pela letra $\pi$ e é dado por $\pi = 1-\beta$.
Uma vez que $\alpha$ está fixado, geralmente podemos aumentar o poder do teste tomando um número maior de amostras.

Agora que já falamos de hipóteses nula e alternativa, do significado de rejeitar ou não rejeitar $H_0$, da diferença entre os erros de tipo I e tipo II, do nível de significância $\alpha$ e do poder $\pi$ de um teste, vamos introduzir mais alguns conceitos, novamente em forma de exemplo.

\begin{example}
\label{exp:peixes}
Sabe-se que a proporção de machos em uma determinada espécie de peixes é de 42\%. Uma bióloga desconfia que, devido a uma alteração na composição bioquímica de um lago, essa proporção é maior.
Assim sendo, temos
\[
H_0:p=0,42
\qquad
H_1:p>0,42
.
\]
Para testar sua conjectura, a bióloga decidiu que iria capturar $n=200$ peixes, observar o sexo deles, calcular a proporção $\hat{p}$ de machos nessa amostra, e comparar com algum nível de corte pré-estabelecido.

Com um nível de significância de 5\%, ela calculou o ponto de corte, que ela chamou ``valor crítico'', como sendo $p_c=0,477$, e decidiu que iria rejeitar a hipótese nula caso $\hat{p} \geq 0,477$.
Para encontrar esse valor de $p_c$, ela usou o fato de que
\[
\hat{p}
\approx
p + Z \cdot \sqrt{\frac{p(1-p)}{n}}
\]
onde $Z$ representa uma normal padrão.
Como $P(Z\geq z_c)=0,05$ para $z_c = 1,645$, ela substituiu o valor de $z$ na igualdade acima, obtendo
\[
p_c
=
p + z_c \cdot \sqrt{\frac{p(1-p)}{n}}
=
0,477
.
\]

Assim, o critério de decisão ficou sendo o seguinte: fixou-se previamente o intervalo $[0,477; 1]$, ao que a bióloga chamou de ``região crítica'', e, após fazer a amostragem, deveria-se rejeitar a hipótese de que $p=0,42$ no caso em que $\hat{p}$ estivesse nesse intervalo.

Perguntando-se sobre o poder desse teste, ela percebeu que não teria como calculá-lo, pois não sabia qual o valor de $p$ deveria considerar, dentre todas as possibilidades que compõem a hipótese alternativa. No máximo, poderia esboçar o gráfico de como seria o poder $\pi$ para distintos de $p$, obtendo a figura abaixo:

\hfil \includegraphics[width=.5\textwidth]{poder1}
\embedfile{figures/poder1.pdf}

Para calcular o valor de $\pi(p)$, ela usou a fórmula
\[
\pi(p) = P_p(\hat{p} \geq p_c) \approx P\Big(Z \geq \frac{p_c-p}{\sqrt{\frac{p(1-p)}{n}}}\Big)
\]
e com isso construiu o gráfico acima.
Os valores de $p \leq 0,42$ que aparecem no gráfico não fazem parte da hipótese alternativa.
Quando $p$ se aproxima de $0,42$, o poder do teste se aproxima de $\alpha$, como indicado no gráfico acima.
Isso corresponde ao fato de que $\alpha$ é justamente a probabilidade de que se rejeite a hipótese nula no caso em que $p=0,42$.

Ao fazer a amostra, a bióloga contou $93$ machos dentre os $200$ peixes amostrados.
Assim, calculou $\hat{p}_{\mathrm{obs}}=0,465$.
Como $0,465<0,477$, ela observou que, supondo a hipótese nula $p=0,42$ como verdadeira, a diferença entre o valor obtido $\hat{p}_{\mathrm{obs}}=0,465$ e o valor teórico $p=0,42$ poderia ser explicada por flutuações aleatórias inerentes ao processo de amostragem, com um nível de significância de $5\%$.
Assim, concluiu que aquela amostra não havia fornecido suficiente evidência estatística para rejeitar a hipótese de que $p=0,42$.
Ela continuou suspeitando que $p>0,42$, mas não conseguiu provar suas suspeitas.
\end{example}

O exemplo acima formalizou um conceito importante, que já aparecia implicitamente nos exemplos anteriores.
A \textbf{região crítica} estabelece o critério para decidir se a amostra oferece evidência estatística suficiente para rejeitar $H_0$.
Se a estatística $\hat{p}$ é tal que $\hat{p} \in RC$, então rejeitamos a hipótese nula. Caso contrário, não rejeitamos a hipótese nula.
Dito de outra forma, quando $\hat{p} \in RC$, dizemos que há evidência estatística para rejeitar $H_0$.

O teste feito nos exemplos acima era do tipo unilateral.
Nesses testes, temos um único valor de $p_c$ e comparamos a estatística observada $\hat{p}$ com esse valor, para decidir se $H_0$ deve ser rejeitada ou não.
Ou seja, a região crítica consistia em um único intervalo que começava em $p_c$ ou terminava em $p_c$.
Mais formalmente, dizemos que o teste é \textbf{unilateral} se $RC$ é da forma $$RC=(-\infty, p_c]$$ ou $$RC=[p_c, +\infty).$$
Do contrário, dizemos que o teste é \textbf{bilateral} se $RC$ é da forma $$RC=(-\infty, p_{c_1}] \cup [p_{c_2}, +\infty),$$
com dois pontos críticos, um para rejeitar a hipótese nula ``pela esquerda'' e outro para rejeitá-la ``pela direita''.
A formulação acima inclui intervalos semi-infinitos porque ela será usada em contextos mais gerais, apesar de que esta seção trata um caso em que a estatística de interesse está sempre em $[0,1]$.

Uma outra observação importante é que a região crítica é construída supondo que $H_0$ é verdadeira.
Apenas sua forma (unilateral, bilateral, etc.)\ leva em conta $H_1$.

Finalmente, passemos a falar da grande novidade do exemplo acima:
a possibilidade de pôr à prova uma hipótese nula que fixa o valor do parâmetro $p$, propondo como alternativa algo mais vago, que não fixa o valor de $p$ e apenas diz que ele parece ser maior, menor, ou diferente do que aquele dado pela hipótese nula.

De forma mais geral, classificamos as hipóteses em dois tipos.
Uma hipótese pode ser \textbf{simples}, se fornece uma distribuição exata para a variável em questão (como por exemplo $p=0,6$ ou $p=0,42$), ou \textbf{composta}, se é dada por uma coleção de distribuições (como por exemplo $p>0,42$).
No caso de uma hipótese alternativa composta, não podemos falar no poder do teste como sendo um número, pois para calcular o poder é necessário especificar qual dos parâmetros incluídos em $H_1$ está sendo considerado.
Nesse caso, consideramos a \textbf{função poder} $\pi$, que associa a cada possível valor do parâmetro populacional $p$ um valor $\pi(p)$, que é o poder do teste no caso de $p$ ser o verdadeiro parâmetro.
A hipótese nula também poderia ser composta, mas essa situação está fora dos objetivos desta seção.

Concluímos esta seção com um último exemplo.

\begin{example}
\label{exp:moeda}
Suspeita-se que uma moeda não é equilibrada, isto é, não tem a mesma probabilidade de sair cara e coroa.
Nesse caso, temos:
\[
H_0:p=\tfrac{1}{2}
\qquad
H_1:p\ne\tfrac{1}{2}
.
\]
Para pôr à prova a hipótese de que a moeda é equilibrada, fixamos o nível de significância $\alpha=0,01$, e lançamos a moeda $n=4000$ vezes.
Antes de lançar a moeda, vamos estabelecer a região crítica.
Primeiro encontramos valores de referência para $z_c$ e depois convertemos para $p_c$ como no exemplo anterior.
Nesse processo, obtemos
\[
P(Z \leq z_{c_1})=\tfrac{\alpha}{2}
\implies
z_{c_1} = -2,58
\implies
p_{c_1} = 0,4796
\]
e
\[
P(Z \geq z_{c_2})=\tfrac{\alpha}{2}
\implies
z_{c_2} = 2,58
\implies
p_{c_2} = 0,5204
.
\]
Assim, estabelecemos a região crítica
\[
RC=
(-\infty; 0,4796] \cup [0,5204; +\infty)
.
\]

O gráfico da função poder ficou assim:

\hfil \includegraphics[width=.5\textwidth]{poder2}
\embedfile{figures/poder2.pdf}

Imaginemos que, após lançar a moeda 4000 vezes, tenhamos observado 1915 coroas.
Com essa observação, calculamos $\hat{p}_{\mathrm{obs}}=0,4788$.
Como $\hat{p}_{\mathrm{obs}} \in RC$, observamos que, supondo a hipótese nula $p=\frac{1}{2}$ como verdadeira, a diferença entre $\hat{p}_{\mathrm{obs}}$ e o valor teórico $p=\frac{1}{2}$ foi grande demais para ser explicada pelas flutuações aleatórias associadas ao processo de lançamento da moeda, com um nível de significância de $1\%$.
Assim, afirmamos haver encontrado suficiente evidência estatística para rejeitar a hipótese de que $p=\frac{1}{2}$.
\end{example}

\subsection*{Cálculos e gráficos}

Abaixo incluímos códigos em R referentes aos cálculos feitos nos exemplos desta seção.
Esses cálculos podem ser reproduzidos com uma calculadora comum e uma tabela normal padrão.
	
\begin{figure}[H]
Exemplo~\ref{exp:facil}
--
Calculando $\alpha$ e $\beta$ para $n=1000$ e $p_c=0,65$:
\footnotesize
\begin{verbatim}
n  <- 1000
pc <- 0.65
p  <- 0.7
zc <- (pc-p)/sqrt(p*(1-p)/n)
zc
# -3.450328
pnorm(zc)
# 0.0002799531
p  <- 0.6
zc <- (pc-p)/sqrt(p*(1-p)/n)
zc
# 3.227486
1 - pnorm(zc)
# 0.0006244155
\end{verbatim}
\end{figure}

\begin{figure}[H]
Exemplo~\ref{exp:dificil}
--
Calculando $\alpha$ para $n=100$ e $p_c=0,65$:
\footnotesize
\begin{verbatim}
n  <- 100
p  <- 0.7
pc <- 0.65
zc <- (pc-p)/sqrt(p*(1-p)/n)
zc
# -1.091089
pnorm(zc)
# 0.1376168
\end{verbatim}
\end{figure}

\begin{figure}[H]
Calculando $p_c$ para $n=100$ e $\alpha=0,001$, calculando $\beta$:
\footnotesize
\begin{verbatim}
n  <- 100
p  <- 0.7
alpha <- 0.001
zc <- qnorm(alpha)
pc <- p+zc*sqrt(p*(1-p)/n)
pc
# 0.5583878
p  <- 0.6
zc <- (pc-p)/sqrt(p*(1-p)/n)
zc
# -0.8494062
1 - pnorm(zc)
# 0.8021723
\end{verbatim}
\end{figure}

\begin{figure}[H]
Calculando $p_c$ para $n=100$ e $\alpha=0,05$, calculando $\beta$:
\footnotesize
\begin{verbatim}
n  <- 100
p  <- 0.7
alpha <- 0.05
zc <- qnorm(alpha)
pc <- p+zc*sqrt(p*(1-p)/n)
pc
# 0.6246233
p  <- 0.6
zc <- (pc-p)/sqrt(p*(1-p)/n)
zc
# 0.5026218
1 - pnorm(zc)
# 0.3076151
\end{verbatim}
\end{figure}

\begin{figure}[H]
Exemplo~\ref{exp:peixes}
--
Calculando $p_c$ para $n=200$, $p=0,42$ e $\alpha=0,05$:
\footnotesize
\begin{verbatim}
n  <- 200
p  <- 0.42
alpha <- 0.05
zc <- qnorm(1-alpha)
pc <- p + zc*sqrt(p*(1-p)/n)
pc
# 0.4774052
\end{verbatim}
\end{figure}

\begin{figure}[H]
Gerando o gráfico da função poder:
\footnotesize
\begin{verbatim}
pdf('poder1.pdf')
p <- seq(.35,.6,.005)
poder <- 1 - pnorm((pc-p)/sqrt(p*(1-p)/n))
plot(p,poder,"l")
lines(c(0.35,.42,.42),c(alpha,alpha,0),lty="dotted")
points(.42,alpha)
dev.off()
\end{verbatim}
\end{figure}

\begin{figure}[H]
Exemplo~\ref{exp:moeda}
--
Calculando $p_{c_1}$ e $p_{c_2}$ para $n=4000$, $p=\frac{1}{2}$ e $\alpha=0,01$.
\footnotesize
\begin{verbatim}
n <- 4000
alpha <- 0.01
p <- 0.5
zc1 <- qnorm(alpha/2)
zc1
# -2.575829
pc1 <- p + zc1*sqrt(p*(1-p)/n)
pc1
# 0.4796363
zc2 <- qnorm(1-alpha/2)
zc2
# 2.575829
pc2 <- p + zc2*sqrt(p*(1-p)/n)
pc2
# 0.5203637
\end{verbatim}
\end{figure}

\begin{figure}[H]
Gerando o gráfico da função poder:
\footnotesize
\begin{verbatim}
pdf('poder2.pdf')
p <- seq(.45,.55,.001)
poder <- 1 - pnorm((pc2-p)/sqrt(p*(1-p)/n)) + pnorm((pc1-p)/sqrt(p*(1-p)/n))
plot(p,poder,"l")
points(.5,alpha)
dev.off()
\end{verbatim}
\end{figure}

\section{Estatística $Z$}

Mais adiante veremos outros tipos de testes de hipótese, que serão cada vez mais complexos.
Na maioria das situações, o teste é feito em termos de uma \emph{estatística}, que não necessariamente é dada pelo \emph{estimador} do parâmetro que está sendo testado.

Em todos os exemplos da seção anterior, a hipótese nula era $H_0:p=p_0$, e o procedimento foi o mesmo.
Fizemos uma Amostra Aleatória Simples de tamanho $n$, que resulta na observação das variáveis $X_1,\dots,X_n$, independentes e com a mesma distribuição, Bernoulli de parâmetro $p$, sendo que a variável $X_k$ indica se a $k$-ésima amostra tem a propriedade estudada (se é uma bola amostrada é de cor viva, se o peixe capturado é macho, se a moeda lançada saiu coroa, etc.).

A partir da amostra $X_1,\dots,X_n$, calculamos o estimador
\[
\hat{p} = \frac{X_1 + \dots + X_n}{n}
.
\]
O resultado do teste é o de rejeitar a hipótese nula se o valor obseravdo do estimador está na \emph{região crítica}.
A região crítica é calculada de forma a satisfazer
\[
P_{H_0}(\hat{p} \in RC) = \alpha.
\]

A \emph{estatística padronizada} para esse teste é definida como
\[
Z = \frac{\hat{p}-p_0}{\ \sqrt{\frac{p_0(1-p_0)}{n}}\ }
.
\]
Supondo $H_0$ válida, essa estatística mede o tamanho do desvio de $\hat{p}$ em torno do seu valor esperado, medido em múltiplos do seu desvio-padrão.
Ainda supondo $H_0$ válida, essa estatística tem distribuição bem aproximada por uma normal padrão, desde que $n$ seja suficientemente grande.
Para os fins práticos de testes de hipótese essa aproximação é razoável quando $n \cdot p_0 \cdot (1-p_0) \geq 5$.
A região crítica para $Z$ pode ser obtida diretamente da tabela normal, sem necessidade de fazer contas envolvendo $n$ e $p_0$.

Para o Exemplo~\ref{exp:dificil}, como a hipótese alternativa era da forma $H_1:p<p_0$, o teste na estatística $Z$ tem como região crítica $(-\infty, z_c]$, onde $z_c$ é tal que $P_{H_0}(Z \leq z_c) = \alpha$.
Como $\alpha=0,05$, da tabela normal obtemos $z_c = -1,645$.
Esse exemplo terminou em suspense, pois não disse o que ocorreu depois que a amostra foi feita.

Para o Exemplo~\ref{exp:peixes}, como a hipótese alternativa era da forma $H_1:p>p_0$, o teste na estatística $Z$ tem como região crítica $[z_c, +\infty)$, onde $z_c$ é tal que $P_{H_0}(Z \geq z_c) = \alpha$.
Como $\alpha=0,05$, da tabela normal obtemos $z_c = 1,645$.
Depois de coletada a amostra, observamos $\hat{p}_{\mathrm{obs}} = \frac{93}{200} = 0,465$.
Calculando a estatística $Z$, obtemos
\[
z_{\mathrm{obs}} = \frac{\hat{p}_{\mathrm{obs}}-p_0}{\ \sqrt{\frac{p_0(1-p_0)}{n}}\ }
=
1,289
.
\]
Como $z_{\mathrm{obs}} \not\in RC$, não rejeitamos a hipótese nula, ou seja, concluímos que não há suficiente evidência estatística indicando que proporção de machos tenha aumentado.

Finalmente, para o Exemplo~\ref{exp:moeda}, como a hipótese alternativa era da forma $H_1:p \ne p_0$, o teste na estatística $Z$ tem como região crítica $(-\infty, -z_{c}] \cup [z_c, +\infty)$, onde $z_c$ é tal que $P_{H_0}(|Z| \geq z_c) = \alpha$.
Como $\alpha=0,01$, da tabela normal obtemos $z_c = 2,576$.
Depois de coletada a amostra, observamos $\hat{p}_{\mathrm{obs}} = \frac{1915}{4000} = 0,47875$.
Calculando a estatística $Z$, obtemos
\[
z_{\mathrm{obs}} = \frac{\hat{p}_{\mathrm{obs}}-p_0}{\ \sqrt{\frac{p_0(1-p_0)}{n}}\ }
=
-2,6879
.
\]
Como $z_{\mathrm{obs}} \in RC$, rejeitamos a hipótese nula, ou seja, concluímos que há suficiente evidência estatística indicando que a moeda não é equilibrada.

\begin{figure}[H]
Código em R para os cálculos desta seção:
\footnotesize
\begin{verbatim}
qnorm(1 - 0.05)
# 1.644854
qnorm(1 - 0.01/2)
# 2.575829
\end{verbatim}
\end{figure}


\section{Nível descritivo ou $p$-valor}

Revisitemos agora o Exemplo~\ref{exp:peixes}. Neste exemplo, a bióloga tinha conjecturado, baseada em estudos teóricos, que poderia haver um aumento na proporção de machos em uma espécie de peixes, devido a uma alteração na composição bioquímica de um lago. Após coletar e analisar 200 amostras, não conseguiu produzir suficiente evidência estatística para apoiar sua tese, com um nível de significância $\alpha=0,05$.

Coletar as 200 amostrar para concluir que não forneciam evidência estatística suficiente foi muito frustrante.
Essa bióloga pode agora se perguntar: o que teria ocorrido se o nível de significância tivesse sido outro?

Para $\alpha=0,02$, certamente não rejeitaríamos $H_0$, já que não foi rejeitado com $\alpha=0,05$, que é maior.
Com $\alpha=0,15$, após consultar a tabela normal e realizar os devidos cálculos, teríamos $p_c = 0,456$ e a hipótese nula teria sido rejeitada.
Logo, com qualquer $\alpha$ maior ou igual a $0,15$, a hipótese nula teria sido rejeitada.

Com $\alpha=0,07$, após consultar a tabela normal e realizar os devidos cálculos, teríamos $p_c = 0,4715$ e a hipótese nula não teria sido rejeitada.
Logo, com qualquer $\alpha$ menor ou igual a $0,07$, a hipótese nula não teria sido rejeitada.

Continuando a busca pelo menor valor de $\alpha$ que levaria à rejeição da hipótese nula, tentamos $\alpha=0,11$.
Após consultar a tabela normal e realizar os devidos cálculos, teríamos $p_c = 0,4628$ e a hipótese nula teria sido rejeitada.
Com $\alpha=0,09$, após consultar a tabela normal e realizar os devidos cálculos, teríamos $p_c = 0,4668$ e a hipótese nula não teria sido rejeitada.
Logo, com qualquer $\alpha$ menor ou igual a $0,09$, a hipótese nula não teria sido rejeitada, e com qualquer $\alpha$ maior ou igual a $0,11$ a hipótese nula teria sido rejeitada.

Se observamos bem, estamos buscando o menor valor de $\alpha$ que levaria à rejeição da hipótese nula usando o estimador $\hat{p}$.
Poderíamos continuar indefinidamente com essa busca, trabalhosa e tediosa, mas podemos obter um resultado imediato a partir da estatística $Z$.
Como $z_{\mathrm{obs}}=1,289$, podemos buscar $P(Z \geq z_{\mathrm{obs}})$ diretamente na tabela normal, sem realizar nenhum novo cálculo, obtendo assim $\hat{\alpha}=0,0987$.
Esse valor é chamado de \emph{nível descritivo} da amostra, também conhecido com o $p$-valor.

Ao invés de fixar um nível de significância $\alpha$ previamente à coleta da amostra, coletar a amostra primeiro, a bióloga pode calcular o $p$-valor $\hat{\alpha}$, e publicar seu relatório apenas com esse valor. Ficaria, então, a cargo do leitor decidir se uma amostra com esse $p$-valor é suficiente evidência estatística para rejeitar a hipótese nula de que a proporção de machos não havia sido alterada.

Formalmente, o nível descritivo é definido como \emph{a probabilidade de que a amostra tinha de fornecer valores pelo menos tão extremos quanto aquele efetivamente observado}.

Se a hipótese alternativa é do tipo $H_1:p>p_0$, o $p$-valor é definido como
\[
\hat{\alpha}
=
P(Z \geq z_{\mathrm{obs}})
.
\]
Se a hipótese alternativa é do tipo $H_1:p < p_0$, o $p$-valor é definido como
\[
\hat{\alpha}
=
P(Z \leq z_{\mathrm{obs}})
.
\]
Se a hipótese alternativa é do tipo $H_1:p \ne p_0$, o $p$-valor é definido como
\[
\hat{\alpha}
=
P(|Z| \geq |z_{\mathrm{obs}}|)
=
2\, P(Z \geq |z_{\mathrm{obs}}|)
.
\]

No caso do Exemplo~\ref{exp:moeda}, como $z_{\mathrm{obs}}=-2,6879$, temos
\[
\hat{\alpha}_{\mathrm{obs}}
=
2\, P(Z \geq 2,6879)
=
0,0072
.
\]

Quanto menor o $p$-valor observado, mais forte a evidência estatística contra a hipótese nula.

\begin{figure}[H]
Código em R para os cálculos desta seção:
\footnotesize
\begin{verbatim}
n  <- 200; p  <- 0.42; p + qnorm( 1 - 0.09 ) * sqrt(p*(1-p)/n)
# 0.4667922
1 - pnorm(1.289)
# 0.09869904
2 * ( 1 - pnorm(2.6879) )
# 0.007190293
\end{verbatim}
\end{figure}


\section{Teste de hipótese para a média}

Trabalho em andamento...


\section{Estatística $T$}

Trabalho em andamento...


\end{document}
